
\documentclass{beamer}
\usefonttheme{serif} % default family is serif
\usepackage{graphicx}
\usepackage{xcolor}
\usepackage{colortbl}
\usepackage{multirow,tabularx}
\usepackage{tikz}
\usepackage{hyperref}
\usepackage{physics}
\usepackage{siunitx}

%% For the name of the chemical compounds
%\usepackage{mhchem}

% Code listing
\usepackage{listings}
\usepackage{xcolor}

% Define Bash syntax highlighting
\lstdefinelanguage{python}{
  morekeywords={if, then, else, fi, for, while, do, done, exit, echo, function},
  morecomment=[l]{\#},
  morestring=[b]",
  sensitive=true
}

% Define listing style
\lstset{
  language=python,
  basicstyle=\ttfamily\small,
  keywordstyle=\color{blue}\bfseries,
  commentstyle=\color{gray},
  stringstyle=\color{red},
  numbers=left,
  numberstyle=\tiny\color{gray},
  stepnumber=1,
  numbersep=5pt,
  breaklines=true,
  breakatwhitespace=true,
  frame=single,
  backgroundcolor=\color{lightgray!20},
  tabsize=2
}

\usepackage{amsthm}
\usepackage{amsmath}
%\usepackage{multilined}

%\usefonttheme{professionalfonts} % using non standard fonts for beamer
%\usefonttheme{serif} % default family is serif
%\usepackage{fontspec}
%\setmainfont{Times New Roman}

% Redefine example environment with separate numbering
\newtheorem{examplex}{Example}
\renewenvironment{example}[1][]{%
  \setbeamercolor{block title example}{use=structure,fg=white,bg=structure.fg!75!black}%
  \setbeamercolor{block body example}{parent=normal text,use=block title,bg=block title.bg!10!bg}%
  \begin{examplex}[#1]%
}{%
  \end{examplex}%
}
% For legend of equations
\newenvironment{conditions*}
  {\par{\abovedisplayskip}\noindent
   \tabularx{\columnwidth}{>{$}l<{$} @{\ : } >{\raggedright\arraybackslash}X}}
  {\endtabularx\par\vspace{\belowdisplayskip}}

% Definition block formatting remains the same
\setbeamertemplate{theorems}[numbered]
\setbeamercolor{block title}{use=structure,fg=white,bg=structure.fg!75!black}
\setbeamercolor{block body}{parent=normal text,use=block title,bg=block title.bg!10!bg}

% Package for figure captions
\usepackage[font=small,skip=2pt]{caption}
% Decrease space/padding after figure captions
\setlength{\belowcaptionskip}{-10pt}

% Define blue color
\definecolor{purple}{rgb}{0.3, 0.0, 0.557}
\definecolor{blue}{rgb}{0.0234375, 0.11328125, 0.609375}

\usepackage{multicol}
\usepackage{wrapfig}

% Use blue theme for the presentation
\setbeamercolor{frametitle}{fg=blue}
\setbeamercolor{title}{fg=blue}
\setbeamercolor{structure}{fg=blue}

\newcommand{\circumeq}{\mathrel{\widehat{=}}}
\newcommand\scalemath[2]{\scalebox{#1}{\mbox{\ensuremath{\displaystyle #2}}}}

% Define affiliations
\newcommand{\affiliation}[2]{#1\textsuperscript{\textcolor{blue}{#2}}}

\title{\textcolor{blue}{Introduction to Linear Algebra}}
\subtitle{\textcolor{blue}{Wigner Summer Camp \\ Data and Compute Intensive Sciences Research Group}}
\author{\textcolor{blue}{Balázs, Paszkál, Vince, Levente, Antal \\ Éva, Hajni}}

\date{\textcolor{blue}{7-11 July 2025}}


% Redefine footline to show only the slide number on the bottom left in blue
\setbeamertemplate{footline}
{
  \leavevmode%
  \hbox{%
  \begin{beamercolorbox}[wd=.1\paperwidth,ht=2.25ex,dp=1ex,left]{page number in head/foot}%
    \hspace{1em} \textcolor{blue}{\insertframenumber}%
  \end{beamercolorbox}}%Wigner Summer Camp

  \vskip0pt%
}


\begin{document}

\begin{frame}
  \titlepage
  \begin{columns}
    \centering
    \column{0.3\textwidth}
    %\includegraphics[width=0.7\textwidth]{img/logo.png}
    \column{0.3\textwidth}
    \centering
    \includegraphics[width=0.8\textwidth]{img/logo.png}
    \column{0.3\textwidth}
    %\includegraphics[width=0.7\textwidth]{img/logo.png}
    \centering
  \end{columns}
\end{frame}

\begin{frame}{Introduction}
    Linear algebra deals with:
    \begin{itemize}
        \item Vectors.
        \item Matrices. (The linear transformations of vector spaces.)
    \end{itemize}

    It is used in many places:
    \begin{itemize}
        \item Physics (velocity, force, ...).
        \item Artificial intelligence - neural networks.
        \item \textbf{Adaptive Law-Based Transformation (ALT).} 
        \item ...
    \end{itemize}
\end{frame}

\begin{frame}{What is a Vector?}
  \begin{itemize}
    \item \textbf{Mathematician's view:} A set that satisfies the 8 axiom.
    \item \textbf{Physicist's view:} An arrow in space with direction and magnitude.
    \item \textbf{Computer Scientist's view:} A 1D array or list of numbers.
    \begin{equation}
    \vec{v} = \begin{bmatrix} 3 \\ -1 \\ 2 \end{bmatrix}, \quad
    \mathbf{w} = \langle 4, 0, -5 \rangle, \quad
    \vec{u} = (1, 2, 3).
    \end{equation}
    \item \textbf{Notation:}
    \begin{itemize}
      \item Boldface: $\mathbf{v}$ (common in CS and mathematical books)
      \item Arrow: $\vec{v}$ (common in physics)
      \item Angled brackets or parentheses
    \end{itemize}
    \begin{example}[Vectors]
        \begin{equation}
        \vec{a} = \begin{bmatrix} 1 \\ 0 \end{bmatrix}, \quad
        \vec{b} = \begin{bmatrix} 1 \\ 1 \end{bmatrix}, \quad
        \vec{c} = \begin{bmatrix} -3 \\ 2 \\ 5 \end{bmatrix}, \quad
        \vec{0} = \begin{bmatrix} 0 \\ 0 \end{bmatrix}.
        \end{equation}
    \end{example}
  \end{itemize}
\end{frame}

\begin{frame}{What is the Dot Product?}
  \begin{itemize}
    \item The dot product of two vectors $\vec{a}$ and $\vec{b}$ is:
    \begin{equation}
    \vec{a} \cdot \vec{b} = a_1 b_1 + a_2 b_2 + \dots + a_n b_n.
    \end{equation}
    \item It can be shown that:
    \begin{equation}
        \vec{a} \cdot \vec{b} = ||\vec{a}||\cdot||\vec{b}||\cos(\angle\left(\vec{a},\vec{b}\right)).
    \end{equation}
    \begin{example}[Dot product]    
        \begin{equation}
        \begin{bmatrix} 2 \\ 3 \end{bmatrix} \cdot \begin{bmatrix} 4 \\ -1 \end{bmatrix} = 2\cdot4 + 3\cdot(-1) = 8 - 3 = 5.
        \end{equation}
    \end{example}
  \end{itemize}
\end{frame}

\begin{frame}{Dot Product - Practice Exercises}
  \begin{itemize}
    \item Calculate the dot product of the following vector pairs:
    \begin{enumerate}
      \item $\begin{bmatrix} 1 \\ -2 \\ 3 \end{bmatrix} \cdot \begin{bmatrix} 4 \\ 0 \\ -1 \end{bmatrix}.$
      \item $\begin{bmatrix} 2 \\ 5 \end{bmatrix} \cdot \begin{bmatrix} -1 \\ 1 \end{bmatrix}.$
    \end{enumerate}
  \end{itemize}
\end{frame}

\begin{frame}{Dot Product - Solutions}
  \begin{itemize}
    \item \textbf{Solution to (a):}
    \[
    1 \cdot 4 + (-2) \cdot 0 + 3 \cdot (-1) = 4 + 0 - 3 = \boxed{1}.
    \]
    \item \textbf{Solution to (b):}
    \[
    2 \cdot (-1) + 5 \cdot 1 = -2 + 5 = \boxed{3}.
    \]
    \item Interpretation:
    \begin{itemize}
      \item (a) Zero \(\Rightarrow\) the vectors are perpendicular.
      \item (b) Positive \(\Rightarrow\) angle between the vectors is acute.
    \end{itemize}
  \end{itemize}
\end{frame}

\begin{frame}{What is a Matrix?}
  \begin{itemize}
    \item A 2D array of numbers arranged in rows and columns:
    \begin{equation}
    A = \begin{bmatrix} 1 & 2 & 3 \\ 4 & 5 & 6 \end{bmatrix},
    \end{equation}
    where \(A\) is a \(2\cross3\) matrix.
    \item Types:
    \begin{itemize}
      \item \textbf{Square:} same number of rows and columns (e.g., $2\times2$, $3\times3$). Mathematicians like them.
      \item \textbf{Non-square:} different number of rows and columns (e.g., $2\times3$, $3\times2$).
    \end{itemize}
  \end{itemize}
\end{frame}

\begin{frame}{Vector-Matrix Multiplication}
  \begin{itemize}
    \item Each row of the matrix is dotted with the vector:
    \begin{equation}
    A = \begin{bmatrix} 1 & 2 \\ 3 & 4 \\ 5 & 6 \end{bmatrix}, \quad \vec{v} = \begin{bmatrix} 2 \\ 1 \end{bmatrix}.
    \end{equation}
    \item Multiplication:
    \begin{equation}
    A\vec{v} = \begin{bmatrix}
    1\cdot2 + 2\cdot1 \\
    3\cdot2 + 4\cdot1 \\
    5\cdot2 + 6\cdot1
    \end{bmatrix} = \begin{bmatrix} 4 \\ 10 \\ 16 \end{bmatrix}.
    \end{equation}
  \end{itemize}
\end{frame}

\begin{frame}{Vector-Matrix Multiplication - Exercises}
  \begin{itemize}
    \item Multiply the vector $\vec{v}$ with matrix $M$ in the following examples:
    \begin{enumerate}
      \item General case:
      \begin{equation}
      \vec{v} = \begin{bmatrix} 1 \\ 2 \end{bmatrix},\quad
      M = \begin{bmatrix} 3 & 0 & 1 \\ 1 & 4 & 2 \end{bmatrix}.
      \end{equation}
      \item Interesting effect:
      \begin{equation}
      \vec{v} = \begin{bmatrix} 1 \\ 1 \end{bmatrix},\quad
      M_1 = \begin{bmatrix} 2 & 0 \\ 0 & 2 \end{bmatrix}, \quad
      M_2 = \begin{bmatrix} -2 & -2 \\ -2 & -2 \end{bmatrix}.
      \end{equation}
      \item Even more interesting effect:
      \begin{equation}
      \vec{v} = \begin{bmatrix} 1 \\ 1 \end{bmatrix},\quad
      M_3 = \begin{bmatrix} 1 & -1 \\ 2 & -2 \end{bmatrix}, \quad
      M_4 = \begin{bmatrix} 5 & -5 \\ -3 & 3 \end{bmatrix}.
      \end{equation}
    \end{enumerate}
    \item Try to interpret the geometric meaning of each result.
  \end{itemize}
\end{frame}

\begin{frame}{Vector-Matrix Multiplication - Solutions}
  \begin{enumerate}
    \item General case:
    \begin{equation}
    \vec{v} \cdot M = \begin{bmatrix} 1 \\ 2 \end{bmatrix} \cdot
    \begin{bmatrix} 3 & 0 & 1 \\ 1 & 4 & 2 \end{bmatrix}
    = \begin{bmatrix} 5 \\ 8 \\ 5 \end{bmatrix}.
    \end{equation}
    \item Stretching cases:
    \begin{equation}
    \vec{v} \cdot M_1 = \begin{bmatrix} 2 \\ 2 \end{bmatrix}=2\vec{v}, \quad
    \vec{v} \cdot M_2 = \begin{bmatrix} -4 \\ -4 \end{bmatrix} = -4\vec{v}.
    \end{equation}
    \item Nearly null output:
    \begin{equation}
    \vec{v} \cdot M_3 = \begin{bmatrix} 0 \\ 0 \end{bmatrix}=0\vec{v}, \quad
    \vec{v} \cdot M_4 = \begin{bmatrix} 0 \\ 0 \end{bmatrix}=0\vec{v}.
    \end{equation}
  \end{enumerate}
\end{frame}

\begin{frame}{Matrices as linear transformation}

\begin{itemize}
    \item Each matrix denotes a transformation $\vec{v} \mapsto M\vec{v}$. It's always linear.
    \item It can be proven that all linear transformation can be described with a matrix.
    
\end{itemize}
    
\end{frame}

\begin{frame}{Transformation examples}
  \begin{itemize}
    \item Stretching: 
    %\begin{equation}
    $\begin{bmatrix} 2 & 0 \\ 0 & 2 \end{bmatrix}.$ %\vec{v} = %\text{doubles length}
    %\end{equation}
    \item Rotation:
    %\begin{equation}
    %\text{Rotation by } 90^\circ: 
    $\begin{bmatrix} 0 & -1 \\ 1 & 0 \end{bmatrix}.$
    %\end{equation}
    \item Reflection:
    %\begin{equation}
    %\text{Reflection over x-axis: }
    $\begin{bmatrix} 1 & 0 \\ 0 & -1 \end{bmatrix}.$
    %\end{equation}
    \item Projection (flattening onto a line or plane) $\begin{bmatrix} 1 & 0 \\ 0 & 0 \end{bmatrix}.$
  \end{itemize}
\end{frame}

\begin{frame}{Matrix-Matrix Multiplication}
  Procedure is the same as the Vector-Matrix multiplication. If an \(n\cross m\) matrix is multiplied by an \(m\cross k\), the result is an \(n\cross k\) matrix.
  \begin{itemize}
    \item The number of columns in the first matrix must be equal to the number of rows in the second matrix.
  \end{itemize}
  
  \[
  A = \begin{bmatrix}
  a_{11} & a_{12} \\
  a_{21} & a_{22}
  \end{bmatrix},
  \quad
  B = \begin{bmatrix}
  b_{11} & b_{12} \\
  b_{21} & b_{22}
  \end{bmatrix}
  \]

  \[
  C = A \cdot B =
  \begin{bmatrix}
  c_{11} & c_{12} \\
  c_{21} & c_{22}
  \end{bmatrix}
  \]

  \[
  \begin{aligned}
  c_{11} &= a_{11} \cdot b_{11} + a_{12} \cdot b_{21} \\
  c_{12} &= a_{11} \cdot b_{12} + a_{12} \cdot b_{22} \\
  c_{21} &= a_{21} \cdot b_{11} + a_{22} \cdot b_{21} \\
  c_{22} &= a_{21} \cdot b_{12} + a_{22} \cdot b_{22}
  \end{aligned}
  \]
\end{frame}

\begin{frame}{Example exercises for matrix-matrix multiplication}
  \begin{example}
    \begin{equation}
      A = \begin{bmatrix} 1 & 2 \\ 3 & 4 \end{bmatrix}, \quad B = \begin{bmatrix} 2 & 0 \\ 1 & 2 \end{bmatrix}
      \Rightarrow A\cdot B = \begin{bmatrix} 4 & 4 \\ 10 & 8 \end{bmatrix}
    \end{equation}
  \end{example}
  \begin{example}
    \begin{equation}
      A = \begin{bmatrix} 1 & 2 & 3 \\ 4 & 5 & 6 \end{bmatrix}, B = \begin{bmatrix} 1 & 0 \\ 0 & 1 \\ 1 & 1 \end{bmatrix}
      \Rightarrow A\cdot B = \begin{bmatrix} 4 & 5 \\ 10 & 11 \end{bmatrix}
    \end{equation}
  \end{example}
  \begin{example}
    \begin{equation}
      A = \begin{bmatrix} 8 & 4 & 11 & 5 \\ 4 & 4 & 21 & 6 \\ 6 & 9 & 11 & 8 \end{bmatrix}, B = \begin{bmatrix} 3 & 2 \\ 4 & 6 \\ 1 & 4 \\ 7 & 2 \end{bmatrix}
      \Rightarrow A\cdot B = \begin{bmatrix} 86 & 94 \\ 91 & 128 \\ 121 & 126\end{bmatrix}
    \end{equation}
  \end{example}
\end{frame}

\begin{frame}{Eigenvalues and Eigenvectors}
  \begin{itemize}
    \item Notice that 
    \begin{equation}
        \begin{bmatrix} 2 & 0 \\ 0 & 2 \end{bmatrix}\begin{bmatrix} 1 \\ 1 \end{bmatrix} = 2\begin{bmatrix} 1 \\ 1 \end{bmatrix}. 
    \end{equation}
    \item Generally for a square matrix \(A\) if
    \begin{equation}
        A\vec{v} = \lambda \vec{v},
    \end{equation}
    then:
    \begin{itemize}
      \item $\vec{v}$ is an eigenvector and
      \item $\lambda$ is the corresponding eigenvalue.
    \end{itemize}
    \item These vectors don't change direction when multiplied by $A$.
  \end{itemize}
\end{frame}

\begin{frame}{Eigenvalues and Eigenvectors - Properties}
  \begin{itemize}
    \item If $A \vec{v} = \lambda \vec{v}$, then $\vec{v}$ is an eigenvector of $A$ with eigenvalue $\lambda$.
    \item Eigenvectors are defined up to a scalar: if $\vec{v}$ is an eigenvector, so is $c\vec{v}$ for any $c \neq 0$.
    \item The set of all eigenvectors corresponding to a single eigenvalue forms a vector subspace.
    \item A square matrix of size $n \times n$ has at most $n$ eigenvalues (including complex ones and multiplicities).
    \item If all eigenvalues of a matrix are positive, the matrix is positive definite.
    \item Diagonal matrices (matrixes where for all $a_{ij}=0$ where $i\neg j$) have their diagonal elements as eigenvalues, and the standard basis vectors as eigenvectors.
    \item The eigenvalues of a complex Hermitian or real symmetric matrix are real.
  \end{itemize}
\end{frame}

\begin{frame}[fragile]{Finding eigenvalues}
Using a computer:
\begin{lstlisting}
import torch
A = torch.tensor([[2., 0.], [0., 2.]])
eigenvalues, eigenvectors = torch.linalg.eigh(A)
print("Eigenvalues:", eigenvalues)
print("Eigenvectors:\n", eigenvectors)
\end{lstlisting}
\end{frame}

\begin{frame}{Revision}
In this presentation you learned about:
    \begin{itemize}
        \item Vectors.
        \item Matrices.
        \item Dot products.
        \item Vector-Matrix and Matrix-Matrix multiplication.
        \item Eigen-decomposition.
    \end{itemize}
    
\end{frame}


\begin{frame}{Extra: Eigenvalue Problem: More Detail}
  \begin{itemize}
    \item Step 1: Start with $A\vec{v} = \lambda \vec{v}$.
    \item Step 2: Rewrite as $(A - \lambda I)\vec{v} = 0$.
    \item Step 3: Solve $\det(A - \lambda I) = 0$ (this gives a polynomial).
    \item Step 4: Find the roots (eigenvalues), then solve the nullspace for each to get eigenvectors.
  \end{itemize}
  \begin{example}
    \begin{equation}
    A = \begin{bmatrix} 2 & 1 \\ 1 & 2 \end{bmatrix} \Rightarrow
    \det\left(\begin{bmatrix} 2 - \lambda & 1 \\ 1 & 2 - \lambda \end{bmatrix}\right) = (2 - \lambda)^2 - 1 = \lambda^2 - 4\lambda + 3.
    \end{equation}
    \begin{itemize}
        \item Roots: $\lambda = 1, 3$.
        \item For $\lambda = 3$: $(A - 3I)\vec{v} = 0 \Rightarrow \vec{v}_1 = \begin{bmatrix} 1 \\ 1 \end{bmatrix}.$
        \item For $\lambda = 1$: $(A - I)\vec{v} = 0 \Rightarrow \vec{v}_2 = \begin{bmatrix} 1 \\ -1 \end{bmatrix}.$
    \end{itemize}
  \end{example}
\end{frame}

\end{document}
